% from ega3-1-fr.tex
\oldpage[III]{149}

noethériens soit fini, il faut et il suffit qu'il soit propre et ait ses fibres
$f^{-1}(y)$ finies (pour tout $y\in\mathfrak Y$).
\end{proposition}

\begin{proof}
Grâce aux définitions (3.4.1 et 4.8.2), on est aussitôt ramené à la même proposition
pour $f_0$ (notations de (\hyperref[III.4.8.1-fr]{4.8.1})), ce qui n'est autre que
(4.4.2).
\end{proof}

\section{Un théorème d'existence de faisceaux algébriques cohérents}
\label{section:III.5-fr}

\subsection{Énoncé du théorème}
\label{subsection:III.5.1-fr}

\begin{env}[5.1.1]
\label{III.5.1.1-fr}
Soient $A$ un anneau adique noethérien, $\mathfrak J$ un idéal de définition de
$A$, de sorte que $A$ est séparé et complet pour la topologie $\mathfrak J$-adique.
Si $Y=\operatorname{Spec}(A)$, le schéma formel affine $\operatorname{Spf}(A)$
s'identifie au complété $\widehat Y$ de $Y$ le long de la partie fermée
$Y'=\mathrm V(\mathfrak J)$ (I, 10.10.1). Soient $X$ un $Y$-préschéma (usuel) de
type fini, $f:X\to Y$ le morphisme structural~; nous désignerons par $\widehat X$
le complété de $X$ le long de la partie fermée $X'=f^{-1}(Y')$, ou encore le
$\widehat Y$-préschéma formel $X\times_Y\widehat Y$~; par
$\widehat f:\widehat X\to\widehat Y$ le prolongement de $f$ aux complétés~; enfin,
pour tout $\mathcal O_X$-Module cohérent $\mathcal F$, nous noterons
$\widehat{\mathcal F}$ son complété, $\mathcal F_{/X'}$, qui est un
$\mathcal O_{\widehat X}$-Module cohérent.
\end{env}

\begin{proposition}[5.1.2]
\label{III.5.1.2-fr}
Les hypothèses et notations étant celles de
(\hyperref[III.5.1.1-fr]{5.1.1}), soit $\mathcal F$ un $\mathcal O_X$-Module
cohérent dont le support est propre sur $Y$ (II, 5.4.10). Les homomorphismes
canoniques (4.1.4)
\[
\rho_i:\mathrm H^i(X,\mathcal F)\longrightarrow
\mathrm H^i(\widehat X,\widehat{\mathcal F})
\]
sont alors des isomorphismes.
\end{proposition}

\begin{proof}
Comme $\mathrm H^i(X,\mathcal F)$ est un $A$-module de type fini (3.2.4), donc
identique à son séparé complété pour la topologie $\mathfrak J$-préadique
($0_{\mathrm I}$, 7.3.6), la proposition n'est qu'un cas particulier de (4.1.10).
\end{proof}

Rappelons que les isomorphismes canoniques $\rho_i$ commutent aux cobords pour
toute suite exacte
$0\to\mathcal F'\to\mathcal F\to\mathcal F''\to0$ de $\mathcal O_X$-Modules
cohérents (($0$, 12.1.6) et (I, 10.8.9)).

\begin{corollary}[5.1.3]
\label{III.5.1.3-fr}
Soient $\mathcal F$, $\mathcal G$ deux $\mathcal O_X$-Modules cohérents tels que
l'intersection de leurs supports soit propre sur $Y$. Alors l'homomorphisme canonique
\begin{equation}
\operatorname{Hom}_{\mathcal O_X}(\mathcal F,\mathcal G)
\longrightarrow
\operatorname{Hom}_{\mathcal O_{\widehat X}}
(\widehat{\mathcal F},\widehat{\mathcal G})
\tag{5.1.3.1}
\label{III.5.1.3.1.eq-fr}
\end{equation}
qui, à tout homomorphisme $u:\mathcal F\to\mathcal G$, fait correspondre son
complété $\widehat u:\widehat{\mathcal F}\to\widehat{\mathcal G}$, est un
isomorphisme. De plus, lorsque le morphisme $f$ est fermé, pour que $\widehat u$ soit
injectif (resp. surjectif), il faut et il suffit que $u$ le soit.
\end{corollary}

\begin{proof}
La première assertion est un cas particulier de (4.5.3), dû encore au fait que le
premier membre de (\ref{III.5.1.3.1.eq-fr}) est un $A$-module de type fini, donc
identique à son séparé complété. Pour démontrer la seconde, notons en vertu de
(I, 10.8.14) que $\widehat u$ est injectif (resp. surjectif) si et seulement s'il existe
un voisinage de $X'$ dans lequel $u$ soit injectif (resp. surjectif). La conclusion
résulte donc du lemme suivant~:
\end{proof}

